\documentclass[11pt]{article}

\usepackage[T1]{fontenc}
\usepackage{lmodern}
\usepackage{amsmath,amssymb,amsthm}
\usepackage[margin=1in]{geometry}
\usepackage{microtype}
\usepackage{needspace}
\usepackage[hidelinks]{hyperref}
\hypersetup{
  pdftitle={Twelve-critical graphs with asymptotic edge count two-fifths n squared},
  pdfauthor={Qiyuan Gu}
}

\newtheorem{theorem}{Theorem}
\newtheorem{lemma}[theorem]{Lemma}
\newtheorem{proposition}[theorem]{Proposition}
\numberwithin{equation}{section}
\newcommand{\col}{\operatorname{col}}

\title{Twelve-critical graphs with\\
\texorpdfstring{\(\left(\frac25-o(1)\right)n^2\)}{(2/5 - o(1)) n squared} edges}
\author{Qiyuan Gu\\
University of Chicago\\
\href{mailto:phoenix1203@uchicago.edu}{\texttt{phoenix1203@uchicago.edu}}}
\date{}

\begin{document}
\maketitle

\begin{abstract}
Let \(f_k(n)\) denote the maximum number of edges in a \(k\)-critical
graph on \(n\) vertices. We construct a sequence of twelve-critical
graphs \(G_s\), with orders tending to infinity, such that
\(e(G_s)/|V(G_s)|^2\to2/5\). This disproves the asymptotic formula
\(f_{12}(n)\sim3n^2/8\), and hence the general asymptotic conjecture
in Erd\H{o}s Problem 917 even when the chromatic number is divisible
by three. The construction combines \(K_5\)-saturated graphs of
sublinear maximum degree with a coloring module of Pegden. Its
criticality follows from explicit eleven-colorings after the deletion
of each type of edge.
\end{abstract}

\section{Introduction}

All graphs are finite and simple. A graph is \(k\)-critical if it has
chromatic number \(k\) and every proper subgraph is
\((k-1)\)-colorable. Write \(e(G)=|E(G)|\), and let \(f_k(n)\) be
the largest number of edges in a \(k\)-critical graph on \(n\) vertices.
One of the questions in Erd\H{o}s Problem 917 \cite{Erdos917} asks
whether, for fixed \(k\ge6\),
\begin{equation}\label{eq:conjecture}
 f_k(n)\sim\frac12\left(1-\frac1{\lfloor k/3\rfloor}\right)n^2.
\end{equation}
Known constructions exceed this coefficient when \(3\nmid k\);
see the discussion in Luo, Ma, and Yang \cite{LMY}.
For \(3\mid k\), the join of \(k/3\) odd cycles of equal order
gives the coefficient in \eqref{eq:conjecture}. We show that this
coefficient is not optimal for \(k=12\).

\begin{theorem}\label{thm:main}
There is a sequence of twelve-critical graphs \(G_s\) such that
\[
 |V(G_s)|\longrightarrow\infty,
 \qquad
 \frac{e(G_s)}{|V(G_s)|^2}\longrightarrow\frac25.
\]
In particular, \(f_{12}(n)\not\sim3n^2/8\).
\end{theorem}

The conclusion concerns the general formula \eqref{eq:conjecture}.
It does not resolve the separate question whether
\(f_6(n)\sim n^2/4\).

The two ingredients are the coloring module \(U(S,T)\) introduced
by Pegden \cite{Pegden} and the saturated graphs of small maximum
degree constructed by Alon, Erd\H{o}s, Holzman, and Krivelevich
\cite{AEHK}. We use five copies of the module. Between their active
sets, we take the complement of a graph obtained from a
\(K_5\)-saturated graph. The maximum-degree condition forces two
colors in each active set to be absent from the other four active
sets. An eleven-coloring would then require a common additional
color in all five sets, contradicting the absence of \(K_5\) in the
saturated graph. Saturation supplies the colorings needed after
an edge is deleted.

\section{A coloring module}
\label{sec:module}

We first recall two elementary properties of critical graphs.
If \(F\) is \(p\)-critical and \(x\in V(F)\), then \(F\) has a
\(p\)-coloring in which a prescribed color occurs only at \(x\):
color \(F-x\) with the other \(p-1\) colors. Also, the join
\(F_1\vee F_2\) of a \(p\)-critical graph and a \(q\)-critical
graph is \((p+q)\)-critical. Its chromatic number is \(p+q\).
Deleting an edge within a factor saves one color. After deleting
a joining edge \(xy\), use singleton colors at \(x\) and \(y\)
and identify those two colors, keeping the remaining palettes
disjoint. Vertex deletion likewise saves a color in one factor.

Let \(S\) be eleven-critical and \(T\) be ten-critical.
The module \(U(S,T)\) consists of disjoint copies of \(S\) and
\(T\), together with an independent set
\[
 A=V(S)\times V(T).
\]
Each \((x,y)\in A\) is adjacent to \(x\in V(S)\) and
\(y\in V(T)\). There are no edges between \(S\) and \(T\), and
no other edges. We call \(A\) the active set and call
\(\{x\}\times V(T)\) the row indexed by \(x\).

The following properties are the case \(U(11,10)\) of Pegden's
module lemma \cite[Lemma 2.5]{Pegden}. We include the proof for the
graphs used here; it does not require the factors to be triangle-free.

\begin{lemma}\label{lem:module}
Fix a palette of eleven colors.
\begin{enumerate}
\item In every proper coloring of \(U(S,T)\) from this palette,
the active set uses at least three colors.
\item Given \(a_0=(x_0,y_0)\in A\) and distinct colors
\(\alpha,\beta,\gamma\), there is a proper coloring in which
the active set uses only \(\alpha,\beta,\gamma\), and
\(\gamma\) occurs on the active set only at \(a_0\).
\item Given an edge \(e\) of \(U(S,T)\) and distinct colors
\(\alpha,\beta\), there is a proper coloring of \(U(S,T)-e\)
in which the active set uses only \(\alpha,\beta\).
\end{enumerate}
\end{lemma}

\begin{proof}
For (1), suppose the active colors are contained in
\(\{\alpha,\beta\}\). Since \(S\) uses all eleven colors, it
contains a vertex colored \(\alpha\) and a vertex colored
\(\beta\). Every active vertex in the first row must have color
\(\beta\), so no vertex of \(T\) can have color \(\beta\).
The second row similarly excludes \(\alpha\) from \(T\).
This leaves only nine colors for \(T\), a contradiction.

For (2), color \(S\) so that \(\alpha\) occurs only at \(x_0\).
Color \(T\) with the ten colors other than \(\alpha\), so that
\(\beta\) occurs only at \(y_0\). Extend these colorings by
\begin{equation}\label{eq:singleton-coloring}
 \col(x,y)=
 \begin{cases}
 \alpha,&x\ne x_0,\\
 \beta,&x=x_0,\ y\ne y_0,\\
 \gamma,&(x,y)=(x_0,y_0).
 \end{cases}
\end{equation}
Each active vertex differs in color from both its neighbors.

For (3), there are four possibilities for the deleted edge.
If \(e\in E(S)\), color both \(S-e\) and \(T\) with the ten
colors other than \(\alpha\), and color all of \(A\) with
\(\alpha\). If \(e\in E(T)\), color \(T-e\) with the nine
colors other than \(\alpha,\beta\). Choose \(x_0\in V(S)\)
and color \(S\) with \(\alpha\) only at \(x_0\). Color the
row indexed by \(x_0\) with \(\beta\) and all other rows with
\(\alpha\).

Suppose instead that \(e\) joins an active vertex
\((x_0,y_0)\) to one of its two structural neighbors.
Use the colorings of \(S\) and \(T\) from the proof of (2).
If the deleted edge joins \((x_0,y_0)\) to \(x_0\), use
\eqref{eq:singleton-coloring} with the color of \((x_0,y_0)\)
changed to \(\alpha\). If it joins \((x_0,y_0)\) to \(y_0\),
change that color to \(\beta\). In either case the only edge
that would become monochromatic is the deleted edge.
\end{proof}

\section{From saturated graphs to twelve-critical graphs}
\label{sec:conversion}

A graph is \(K_5\)-saturated if it contains no \(K_5\), but
adding any missing edge creates a \(K_5\). Thus the common
neighborhood of every pair of distinct nonadjacent vertices
contains a triangle.

\begin{proposition}\label{prop:conversion}
Let \(H\) be a \(K_5\)-saturated graph on \(v\ge5\) vertices,
with maximum degree \(d<v-1\). Let \(h\ge11\) be odd and suppose
\begin{equation}\label{eq:degree-condition}
 v>40hd.
\end{equation}
Put \(a=2hv\). There is a twelve-critical graph \(G\) with
\begin{equation}\label{eq:order}
 |V(G)|=5(a+h+2v)
\end{equation}
and
\begin{equation}\label{eq:edge-bound}
 e(G)\ge10a^2\left(1-\frac d v\right).
\end{equation}
More precisely, if \(m=e(H)\), then
\begin{equation}\label{eq:exact-edges}
 e(G)=10(a^2-8h^2m)+5(2a+9h+16v-79).
\end{equation}
\end{proposition}

\begin{proof}
\emph{Construction.}
Set \(\ell=2v\) and take
\[
 S=K_8\vee C_{h-8},\qquad T=K_7\vee C_{\ell-7}.
\]
Both cycles are odd and have length at least three. Hence \(S\)
and \(T\) are eleven-critical and ten-critical, respectively,
with \(h\) and \(\ell\) vertices.

Take five disjoint copies \(U_1,\ldots,U_5\) of \(U(S,T)\),
with active sets \(A_1,\ldots,A_5\), each of size \(a=h\ell\).
In each copy, identify \(V(T)\) with \(V(H)\times\{0,1\}\)
by any bijection. An active vertex then has the form
\((x,(z,\varepsilon))\); give it label \(z\).
Let \(\pi\) denote this labeling map. Each label occurs
exactly \(b=2h\) times in each active set.

Define a five-partite graph \(Z\) on \(A_1\cup\cdots\cup A_5\)
by putting, for \(u\in A_i\) and \(w\in A_j\) with \(i\ne j\),
\begin{equation}\label{eq:Z}
 uw\in E(Z)\quad\Longleftrightarrow\quad
 \pi(u)\pi(w)\in E(H).
\end{equation}
The graph \(G\) retains all the edges of the five modules and,
between different active sets, takes the complement of \(Z\).
There are no other edges between modules. In particular, active
vertices with equal labels in different parts are adjacent in \(G\).

We record three properties of \(Z\). First, \(Z\) is
\(K_5\)-free: the labeling map sends every clique injectively
to a clique of \(H\). Moreover, adding any missing edge between
different parts of \(Z\) creates a \(K_5\) with one vertex in
each part. To see this, let \(u,w\) be the endpoints of such a
missing edge. If their labels \(z,z'\) differ, saturation of
\(H\) gives a triangle in \(N_H(z)\cap N_H(z')\). Choose
vertices with these three labels in the other three parts.
Together with \(u,w\), they form the required clique in
\(Z+uw\). If the labels coincide, say both equal \(z\), choose
a vertex \(z'\) not adjacent to \(z\); this is possible because
\(d<v-1\). The same common-neighborhood triangle is contained
in \(N_H(z)\) and gives the required clique.

Second, \(Z\) contains a \(K_4\) across any prescribed four
parts. Indeed, \(H\) has a missing edge; one endpoint together
with a triangle in the common neighborhood gives a \(K_4\) in
\(H\). Its four labels can be placed in any four parts of \(Z\).

Third, writing \(D=\Delta(Z)\), condition \eqref{eq:degree-condition} gives
\begin{equation}\label{eq:D}
 D=4bd=8hd,\qquad \ell>10D,
\end{equation}

\smallskip
\noindent\emph{The chromatic lower bound.}
Suppose that \(G\) has an eleven-coloring. In each \(A_i\),
at least two color classes have size greater than \(D\).
Otherwise, choose the unique color with more than \(D\) vertices,
if there is one, and call it \(\alpha\); if there is none,
choose any color. The other ten colors together occupy at most
\(10D<\ell\) vertices of \(A_i\). Every row, having \(\ell\)
vertices, therefore contains a vertex of color \(\alpha\).
Each vertex of the copy of \(S\) is adjacent to its entire row,
so this copy of \(S\) would avoid \(\alpha\), contradicting
\(\chi(S)=11\).

A color used more than \(D\) times in \(A_i\) cannot occur in
any other active set. A vertex of that color in another part
would have to be adjacent in \(Z\) to all those more than
\(D\) vertices, contrary to \(\Delta(Z)=D\).
Choose two such colors in each part. These are ten distinct
colors, each absent from the other four active sets. By
Lemma~\ref{lem:module}(1), every active set needs a third color.
The one remaining color must consequently occur in all five
active sets. Choosing a vertex of this color in each part gives
a \(K_5\) in \(Z\), a contradiction. Thus \(\chi(G)\ge12\).

\smallskip
\noindent\emph{Colorings after edge deletion.}
Use an eleven-color palette
\[
 \alpha_1,\beta_1,\ldots,\alpha_5,\beta_5,\gamma.
\]
The pair \(\alpha_i,\beta_i\) will be used only in \(A_i\)
among the active sets. Structural vertices may use the full
palette, since they have no neighbors in other modules.

Let \(e=uw\) join two different active sets in \(G\).
By the first property of \(Z\), there is a \(K_5\) in
\(Z+uw\) containing \(uw\), with vertices \(a_i\in A_i\).
Apply Lemma~\ref{lem:module}(2) in each module, using active
colors \(\alpha_i,\beta_i,\gamma\) and assigning \(\gamma\)
only to \(a_i\) in \(A_i\). The only color shared by different
active sets is \(\gamma\). The five vertices of this color are
pairwise nonadjacent in \(G-e\), because they form a clique in
\(Z+uw\). We have obtained an eleven-coloring of \(G-e\).

Now let \(e\) lie within \(U_i\). Apply
Lemma~\ref{lem:module}(3) to color \(U_i-e\), using only
\(\alpha_i,\beta_i\) on \(A_i\). Choose a \(K_4\) in \(Z\)
across the other four parts. In each of those four modules,
apply Lemma~\ref{lem:module}(2) with the selected vertex as
the only active vertex of color \(\gamma\). Once again this
is a proper eleven-coloring of \(G-e\).

These cases include every edge of \(G\). Recoloring one endpoint
of a deleted edge with a twelfth color gives \(\chi(G)\le12\).
Thus \(\chi(G)=12\), and deletion of every edge lowers the
chromatic number. The graph has no isolated vertices, so a
subgraph missing a vertex is also contained in some \(G-e\).
Consequently \(G\) is twelve-critical.

\smallskip
\noindent\emph{Counting vertices and edges.}
Each module has \(a+h+2v\) vertices, giving \eqref{eq:order}.
Between a fixed pair of active sets, \(Z\) has \(2b^2m\) edges:
each edge of \(H\) gives two ordered pairs of labels, with
\(b^2\) pairs of active vertices for each. Thus the number of
edges of \(G\) between active sets is
\[
 10(a^2-2b^2m)=10(a^2-8h^2m).
\]
The two structural graphs have
\[
 e(S)=9h-44,\qquad e(T)=16v-35,
\]
and each module has \(2a\) further edges incident with its active
set. This proves \eqref{eq:exact-edges}. Finally, \(2m\le vd\)
and \(a=bv\) give \eqref{eq:edge-bound} from the edges between
active sets alone.
\end{proof}

\Needspace{10\baselineskip}
\section{Choice of the saturated graphs}
\label{sec:parameters}

We use the following instance of the construction in
\cite[Section 4, Example 2]{AEHK}.

\begin{lemma}[Alon--Erd\H{o}s--Holzman--Krivelevich]
\label{lem:saturation}
For every prime power \(Q\ge4\), there is a \(K_5\)-saturated
graph \(H_Q\) with
\begin{equation}\label{eq:H-parameters}
 |V(H_Q)|=13Q^2+12Q,
 \qquad \Delta(H_Q)=22Q-3.
\end{equation}
\end{lemma}

To obtain this graph from Example 2 of \cite{AEHK}, set \(k=5\)
and give each of the \(Q^2\) additional independent sets one
vertex. The core has \(12Q(Q+1)\) vertices. Each core vertex has
\(3(3Q-1)\) neighbors in its level, \(12Q\) in the other levels,
and \(Q\) among the additional vertices. Its degree is therefore
\(22Q-3\). Each additional vertex has degree \(12(Q+1)\),
which is smaller for \(Q\ge4\). This gives
\eqref{eq:H-parameters}; saturation is part of the cited construction.

\begin{proof}[Proof of Theorem~\ref{thm:main}]
For each integer \(s\ge4\), set
\[
 h=3^s,\qquad Q=h^2,
\]
and take \(H=H_Q\) from Lemma~\ref{lem:saturation}. Write
\(v=|V(H)|\) and \(d=\Delta(H)\). Then
\[
 v=13h^4+12h^2,\qquad d=22h^2-3<v-1.
\]
The integer \(h\) is odd and at least \(81\). Moreover,
\[
 40hd<880h^3<13h^4\le v,
\]
since \(13h\ge13\cdot81>880\).
Proposition~\ref{prop:conversion} therefore gives a twelve-critical
graph \(G_s\). Put \(a=2hv\) and \(N_s=|V(G_s)|\). We have
\begin{equation}\label{eq:asymptotics}
 \frac{N_s}{5a}=1+\frac1h+\frac1{2v}\longrightarrow1,
 \qquad \frac d v\longrightarrow0.
\end{equation}
There are at most \(10a^2\) edges between active sets, and the
five modules contain \(5(2a+9h+16v-79)=O(a)\) edges.
Together with \eqref{eq:edge-bound}, this gives
\[
 10a^2\left(1-\frac d v\right)
 \le e(G_s)\le10a^2+O(a).
\]
Dividing by \(N_s^2\) and using \eqref{eq:asymptotics} yields
\[
 \frac{e(G_s)}{N_s^2}\longrightarrow\frac{10}{25}=\frac25.
\]
Finally, \(N_s\to\infty\) and \(f_{12}(N_s)\ge e(G_s)\).
Since \(2/5>3/8\), this subsequence rules out
\(f_{12}(n)\sim3n^2/8\).
\end{proof}

\section*{AI tool disclosure}

OpenAI Codex (GPT-6) was used substantially in the development of
this work, including the composition of saturated graphs with the
coloring module, the proof of criticality, literature retrieval,
and preparation and checking of the manuscript. The author selected
the problem and directed the research. An exact-arithmetic script
was also used to check the parameter counts; the proof is independent
of that computation. The author is responsible for the mathematical
content. This draft is circulated for review, and the AI-assisted
checks are not presented as independent human verification.

\begin{thebibliography}{9}

\bibitem{AEHK}
N.~Alon, P.~Erd\H{o}s, R.~Holzman, and M.~Krivelevich,
On \(k\)-saturated graphs with restrictions on the degrees,
\emph{J. Graph Theory} \textbf{23} (1996), 1--20.
\url{https://web.math.princeton.edu/~nalon/PDFS/aehk1.pdf}.

\bibitem{Erdos917}
T.~F. Bloom,
\emph{Erd\H{o}s Problem \#917},
\url{https://www.erdosproblems.com/917}.
Accessed 5 September 2026.

\bibitem{LMY}
C.~Luo, J.~Ma, and T.~Yang,
On the maximum number of edges in \(k\)-critical graphs,
\href{https://arxiv.org/abs/2301.01656}{arXiv:2301.01656}, 2023.

\bibitem{Pegden}
W.~Pegden,
Critical graphs without triangles: an optimum density construction,
\emph{Combinatorica} \textbf{33} (2013), 495--512.
\href{https://arxiv.org/abs/1101.4417v2}{arXiv:1101.4417v2}.

\end{thebibliography}

\end{document}
